% 这是注释，不会出现在pdf文档里

% ----------------------------------
% 支持中文，使用 ctexart， 编码需要设置为GBK
% 使用ctexart时，\section默认是居住排版的
%-----------------------------------
\documentclass[]{ctexart}	    

% === 设置标题的格式 ===
\ctexset{
	section = {
		format+= \zihao{-4} \heiti \raggedright,	% 控制左对齐
		name={,、},
		number = \chinese{section},
		beforeskip=1.0ex plus 0.2ex minus .2ex,
		afterskip = 1.0ex plus 0.2ex minus .2ex,
		aftername = \hspace{0pt}
	},
	subsection = {
		format+= \zihao{5} \heiti \raggedright,	% 控制左对齐
		name={,、},
		number = \arabic{section},		% 章节编号，阿拉伯数字
		beforeskip=1.0ex plus 0.2ex minus .2ex,
		afterskip = 1.0ex plus 0.2ex minus .2ex,
		aftername = \hspace{0pt}
	}
}

%\documentclass[
%10pt,		% 文字默认字体大小，一般只有10,11,12 pt
%]{article}			% 使用article时，不支持中文

% 编辑中文内容需要引用的宏包
\usepackage{ctex}
% 使用图片
\usepackage{graphicx}

% 不带编号的公式使用，使用矩阵
\usepackage{amsmath}

% 多行数学公式使用
\usepackage{amssymb}

% 使用文件的参考文献， plain： 排版样式 ， 还有 alpha
\bibliographystyle{plain}
% 使用 natbib 宏包可以使用更多的样式
%\usepackage[round]{natbib}

% 指定存放图片的文件夹路径
\graphicspath{{pictures/},{photoes/}}

% 定义新命令
\newcommand{\mycommend}{\textit{\textbf{\textsf{Fancy Text}}}}

% 文章开始
\begin{document} 
% == 1. 中文，使用 ctex包， xelatex 编译模式
	你好，中国
	
% == 2. 数学公式
 	hello $f(x)=3x^2-x-1$
 	$$ f(x)=x^y-z+x$$
 	$$ f(x)=(x-y)^p/(z-x)^1/p $$
 	
% == 3. 字体族设置 （罗马字体，无衬线字体，打印机字体）
 	% 罗马字体
 	\textrm{这些内容用罗马字体显示 Roman Family}
 	
 	% 无衬线字体
 	\textsf{这些内容用无衬线字体显示 Sans Serif Family}
 	
 	% 打印机字体
 	\texttt{这些内容用打印机字体显示 Typewriter Family}
 	
 	% 声明后续字体为罗马字体
 	\rmfamily 声明后续字体都为罗马字体
 	
 	% 声明后续字体都为无衬线字体
 	\sffamily 声明后续字体都为无衬线字体
 	
 	% 声明后续字体都为打印机字体
 	\ttfamily 声明后续字体都为打印机字体
 	
 	% 使用大括号限定字体作用范围
 	{\sffamily 这些内容限定使用无衬线字体}
 	
 	这是后续内容，当遇到其他字体声明时，会启用新的字体声明，否则延续上文声明的字体,
 	
 	% 字体系列设置(粗细、宽度)
 	\textmd{中等尺寸字体,Medium Series}  \\
	\textbf{加粗字体, Boldface Series} 	\\
	% 字体形状设置(直立、斜体、伪斜体、小型大写)
	\textup{直立字体，Upright Shape}		\\
	\textit{斜体，Italic Shape}		\\
	\textsl{伪斜体，Slanted Shape}	\\
	\textsc{小型大写，Small Caps Shape}	\\
	% 设置影响后续内容
	{\upshape 影响后续内容为直立字体，大括号限制作用域 Upright Shape} \\
	{\itshape 影响后续内容为斜体字体，大括号限制作用域 Italic Shape} \\
	{\slshape 影响后续内容为伪斜体字体，大括号限制作用域 Slanted Shape} \\
	{\scshape 影响后续内容为小型大写字体，大括号限制作用域 Small Caps Shape} \\
	% 中文字体
	{\songti 宋体} \quad 
	{\heiti 黑体} \quad
	{\fangsong 仿宋} \quad
	{\kaishu 楷书} \quad 	
	\quad 是空格 \\
	{\textbf 中文粗体使用黑体表示}\\
	{\textit 中文斜体使用楷书表示} \\
	
	% 字体大小，相对于 \documentclass 声明中的 normal 字体大小
	{\tiny 			hello}\\
	{\scriptsize 	hello}\\
	{\footnotesize	hello}\\
	{\small			hello}\\
	{\normalsize	hello}\\
	{\large			hello}\\
	{\Large			hello}\\
	{\LARGE			hello}\\
	{\huge			hello}\\
	{\Huge			hello}\\
	
	% 中文设置字号
	\zihao{-0} 初号字 \\
	\zihao{5} 5号字   \\
	
	% 使用自定义命令
	\mycommend
	
% == 4. 章节结果
	% 产生目录
	\tableofcontents

	% section 构建章节， 加*号去掉自动编号
	\section*{引言}
	\subsection*{小节}
	这是引言的内容\\ 这是换行，首行并不缩进，换行符合不能出行在一行的开头。
	
	使用空行也能换行，一个空行或多个空行效果是一样的，这种换行会首行缩进。
	
	\paragraph{段落标题} 产生新的段落
	
	\section{第一章}
	% subsection 构建小节
	\subsection{第一小节}
	\subsubsection{三级小节}
	
	\section{第二章}
	\subsection{第一小节}
	
	\section{第三章}
	\subsection{第一小节}
	
	% \documentclass 是 texbook 时可以使用 \chapter{} 命令
 	

% == 5. 特殊字符
	\section{空白符}
	% 空行分段：多个空行等同与一个
	% 自动缩进：绝对不能使用空格代替
	% 空格：英文中多个处理为1个，中文中空格被忽略
	% 汉字与其他字符的间距会自动有 XeLaTex 处理
	% 禁止使用中文全角空格
	% 使用空白符时，如下：
	a\quad b	% 1em
	a\qquad b	% 2em
	a\thinspace b 或 a\,b	% 约位 1/6个em
	a\enspace b	% 0.5em
	a~b			% 硬空格
	% 1pc=12pt=4.218em
	a\kern 1pc b
	a\kern -1em b
	a\hskip 1emb
	a\hspace{35pt}b
	
	a\hphantom{xyz}b	%占位宽度
	a\hfill b			% 弹性长度
	
	\section{控制符}
	% & #、$、%、{、}、~、等符号都有特殊意义，要输入原符号时需要加反斜线转义
	\# \$ \{ \} \_{} \~{} \^{} 
	这是一个换行符号：\textbackslash 
	
	\&

	\section{排版特殊符号}
	\S \P \dag \ddag \copyright \pounds
	
	\section{\TeX 标志符号}
	\TeX{}
	\LaTeX{}
	\LaTeXe{}
	
% == 6. 插入图片
	\section{插入图片}
	% 使用 \usepackage{graphicx} 宏包
	% 导言区： \usepackage{graphicx}
	% 语  法： \includegraphics[<选项>]{<文件名>}
	% 文件格式：\ EPS， PDF，PNG, JPEG, BMP
	\includegraphics{imgtest}	% 可以省略文件后缀名
	% 通过可选参数指定缩放、高度
	\includegraphics[
	%width=13cm,		% 图片宽度
	%height=8cm,		% 图片高度
	%width=0.1\textwidth,	% 相对宽度
	%height=0.1\textheight,	% 相对高度
	scale=0.5,			% 图片缩放
	]{imgtest.jpg}
	
% == 7. 表格
	\section{插入表格}
	% 使用 tabular生成表格
	% 生成5列的表格 l:左对齐，c:居中对齐, r:右对齐 , | 产生表格竖线， || 产生双竖线
	% p 指定列宽度
	\begin{tabular}{l || c  c | p{1.5cm} | r}
		\hline   % 产生表格横线	
		姓名 & 语文 & 数学 & 外语 & 备注 \\ % 用 & 作为列分隔符 ，\\ 产生新行	
		\hline \hline  % 产生双横线
		张三 & 100 & 100 & 100 & 优秀 \\	
		\hline   % 产生表格横线
	\end{tabular}

	% 使用 booktab 宏包实现三线表
	% 使用 longtab 宏包实现跨页长表格
	
	
% == 8. 浮动环境
	% 使用浮动体可以：实现灵活分页
	% 给图片添加标题
	% 引用图表编号
	% 排版参数：
	% h : 此处(here) 代码所在的上下文位置
	% t : 页顶(top), 代码所在页或之后页的顶部
	% b : 页底(bottom), 代码所在页或之后页的底部
	% p : 独立一页(page)

	% 图片浮动
	\section{图片浮动}
	想看美女吗，见图\ref{mm}
	\begin{figure}[htbp]
		\centering			% 居中
		
		\includegraphics[
		%width=13cm,		% 图片宽度
		%height=8cm,		% 图片高度
		%width=0.1\textwidth,	% 相对宽度
		%height=0.1\textheight,	% 相对高度
		scale=0.5,			% 图片缩放
		]{imgtest.jpg}
		\caption{美女}		% 设置图片标题，命令在图片下则标题在图片下，自动编号
		\label{mm}			% 生成标签
	\end{figure}
	
	% 表格浮动
	\section{表格浮动}
	想看成绩吗 ，见表 \ref{firsttable}
	\begin{table}[htbp]		% [位置参数]
		\centering			% 居中
		\caption{成绩表}		% 设置表格标题，自动编号
		\begin{tabular}{l || c  c | p{1.5cm} | r}
			\hline   % 产生表格横线	
			姓名 & 语文 & 数学 & 外语 & 备注 \\ % 用 & 作为列分隔符 ，\\ 产生新行	
			\hline \hline  % 产生双横线
			张三 & 100 & 100 & 100 & 浮动表格 \\	
			\hline   % 产生表格横线
		\end{tabular}
		\label{firsttable}
	\end{table}
 	
 	
% == 9. 数学公式
	\section{数学模式—基本数学公式}
	
	\subsection{行内公式}
	% 使用单个美元符号排版
	$ x^2+y^2=1 $
	
	% 使用小括号排版
	\( x^2+y^2=1 \)
	
	% 使用math环境排版
	\begin{math}
		x^2+y^2=1
	\end{math}
 	
 	\subsection{上下标}
 	% 上标使用 ^ 与{}配合，否则只将一位数字认为是上标
 	$ 3x^20 + 2y^10 =0 $ 
 	
 	$ 3x^{20} + 2y^{10} =0 $
 	
 	$ x^{p^2+q^2}+y^{n^2+m^2}=0 $
 	% 下标使用 _ 实现
 	$ a_2 + b_2 =0 $
 	
 	$ a_20 + b_20 =0 $
 	
 	$ a_{20} + b_{20} =0 $
 	
 	\subsection{希腊字母}
 	$ \alpha $		\\
 	$ \beta $		\\
 	$ \gamma $		\\
 	$ \epsilon $	\\
 	$ \pi $			\\
 	$ \omega $		\\
 	$ \Gamma $		\\
 	$ \Delta $		\\
 	$ \Theta $		\\
 	$ \Pi $			\\
 	$ \Omega $		\\
 	
 	$ \alpha^2 + \beta^2 + \gamma^2 =1 $
 	
 	\subsection{数学函数}
 	$ \log $
 	$ \sin $
 	$ \cos $
 	$ \arcsin $
 	$ \arccos $
 	$ \ln $
 	$ \log $
 	
 	$ \sin^2x + \cos^2x = 1 $
 	
 	% 排版根式
 	$ \sqrt{2} $
 	$ \sqrt[3]{9} $
 	$ \sqrt[3]{ \sin^2x + \cos^2y + \sqrt{5} } $
 	
 	\subsection{分式}
 	% 一行
 	$ 3/4 $
 	
 	% 上下结构
 	$ \frac{3}{4} $
 	
 	\subsection{行间公式}
 	% $$ 两个美元符号
 	$$ \cos^2 + \sin^2 = 1 $$ 
 	% 中括号
 	交换律是： \[ a+b = b+a \]
 	% 使用 displaymath
 	\begin{displaymath}
 		\cos^2 + \sin^2 = 1
 	\end{displaymath}
 	
 	\subsection{对公式编号}
 	% 使用 equation
 	数学公式见式 \ref{eq1}	% 引用公式
 	\begin{equation}
 		x^2 + y^2 =1 
 		\label{eq1}
 	\end{equation}
 	
 	\subsection{不对公式编号}
 	不编号的公式，见式\ref{eq2} ，使用章节编号
 	% 使用 equation* 需要使用宏包 amsmath
 	\begin{equation*}
 		x^2 + y^2 =1 
 		\label{eq2}
 	\end{equation*}
 	
% == 10. 排版矩阵
	\section{数学模式—矩阵排版}
	% 使用 matrix ,引入 amsmath 宏包
	% 与tabular 表格一样，用 & 分割列，用 \\ 分割行
	% 在矩阵中也可以使用数学公式

	\[
		\begin{matrix}
		0 & 1 \\
		1 & 0
		\end{matrix}
	\]
	
	% pmatrix 用于在矩阵两端加小括号
	\[
		\begin{pmatrix}
			0 & 1 \\
			1 & 0
		\end{pmatrix}
	\]
	
	% bmatrix 加中括号
	\[
		\begin{bmatrix}
			0 & 1 \\
			1 & 0
		\end{bmatrix}
	\]
	
	% Bmatrix 加大括号
	\[
		\begin{Bmatrix}
			0 & 1 \\
			1 & 0
		\end{Bmatrix}
	\]
	
	% vmatrix 加单竖线
	\[
		\begin{vmatrix}
			0 & 1 \\
			1 & 0
		\end{vmatrix}
	\]
	
	% Vmatrix 加双竖线
	\[
		\begin{Vmatrix}
			0 & 1 \\
			1 & 0
		\end{Vmatrix}
	\]
 	
 	% 矩阵中使用上下标
 	\[
 		\begin{pmatrix}
 		a_{11}^{2} & a_{12}^{2} & a_{13}^{2} \\
 		0	& a_{22} & a_{23} \\
 		0 & 0 & a_{33} \\
 		\end{pmatrix}
 	\]
 	
 	% 常用省略号 
 	\(	\dots  \vdots  \ddots \)
 	
 	% 带省略号的矩阵
 	
 	\[
 		\begin{bmatrix}
 		a_{11} & \dots	& a_{1n} \\
 		       & \ddots & \vdots \\
 		0	   & 		& a_{nn} \\
 		\end{bmatrix} _{n \times n}
 	\]
 	
 	% 数学模式中，可以使用 \times 排版乘号
 	$ x \times y =1 $
 	
 	% \text 用于在数学模式中临时切换到文本模式
 	$ \text{数学公式是：} x \times y = 2 $
 	
 	% smallmatrix 用于排版行内小矩阵
 	我们来看一个行内小矩阵 
 	\begin{math}
 		\left( % 需要手动加括号，小括号、中括号 或竖线
 		\begin{smallmatrix}
 		x & y \\
 		1 & 2 \\
 		\end{smallmatrix}
 		\right)
 	\end{math}
 	
 % == 数学模式多行公式排版
 	\subsection{数学模式-多行公式}
 	% 需要引入 amsmath 、amssymb 宏包
 	% 可以使用gather环境实现多行公式的排版
 	% gather环境默认对公式编号，如果要取消编号，使用 gather*
 	% 也可以在 \\ 前使用 \notag 阻止该行编号
 	\begin{gather}
 		a + b = b + a \notag \\    % 使用 \\ 换行
 		a^2 + b^2 + c^2 =0 
 	\end{gather}
 	
 	% 可以使用 align 和 align* 环境 （用*号进行对齐)
 	% align* 不对公式编号
 	\begin{align}
 		x  &= t & x &=\cos t & x &=t \\
 		y & = 2t & y &=\sin(t+1) & y &= t 
 	\end{align}
 	
 	% \in 输出属于符号， \mathbb 输出花体字符，需要 amssymb 宏包
 	\begin{equation}
 		D(x) = \begin{cases}
 			1, & \text{如果} x \in \mathbb(Q); \\
 			0, & \text{如果} x \in \mathbb(R) \setminus \mathbb{Q}.
 		\end{cases}
 	\end{equation}
 	
 	% 常用数学公式
 	\subsubsection{常用数学公式}
 	% 数学求和公式
 	\[
 		\sum_{i=1}^{N}(x^2 + y^2)^{i}
 	\]
 	
 	% 明科夫斯基距离公式
 	明科夫斯基距离公式:
 	\begin{equation} 		
 		dist(X,Y) = (\sum_{i=1}^{N}|hash(x_i) - hash(y_i)|^p)^{1/p}
 	\end{equation}
 	 	
 	 	
 	 
 % == 参考文献
 	\section{参考文献}
 	\subsection{简单排版参考文献}
 	% 简单的参考文件排版
 	% 参考文献格式：
 	% \begin{thebibliography}{编号样本}
 	%	\bibitem[记号]{引用标志} 文献条目1
 	%		\bibitem[记号]{引用标志} 文献条目1
 	% \end{thebibliography}
 	% 引用参考文献 \cite
 	% \emph 强调文献中的某些内容
 	
 	这句话引用参考文献
 	\begin{thebibliography}{99}
 		\bibitem{article1}陈立辉，苏伟，蔡川. \emph{基于LaTex的Web数学公式提取方法研究}[J]. 据算计科学。i2014(06)
 		\bibitem{article2}陈立辉，苏伟，蔡川. \emph{基于LaTex的Web数学公式提取方法研究}[J]. 据算计科学。2014(06) \texttt{设置字体}
 	\end{thebibliography}
 
 	\subsection{单独排版参考文献}
 	% 在option中将构建操作的默认文献工具设置为 BibTex
 	% 创建新文件编辑参考文献的详细信息
 	1. 引用的参考文献才会在下边展示出来，没有引用的不会展示出来。参考文献会按照引用自动编号
 	这是一个引用了的参考文献 \cite{bookLable}
 	
 	2. 可以在百度学术引用里点击 BibTex 引用，生成引用格式，粘贴到参考文献文件里。这里引用 
 	\cite{韩秋明2018区块链技术国外研究述评}
 	知网需要安装 zotero的firefox扩展包
 	
 	3. 将未引用的参考文献全部列出来 
 	
 	4. 引用样式 \citet{bookLable}
 	
 		
 	\nocite{*}
 	
 	\bibliography{bibliography.bib}
 	 
 % 文章结束
\end{document}

